% ME52037 Design Optimisation - Task 2: Black-box optimisation of MCSI
% Compile with pdflatex (e.g. Overleaf). Figures are in ./figures/.
% Target: 6 pages at 11pt, 2 cm margins.
\documentclass[11pt,a4paper]{article}
\usepackage[margin=2cm]{geometry}
\usepackage{graphicx}
\usepackage{amsmath}
\usepackage{booktabs}
\usepackage{xcolor}
\usepackage[font=small,labelfont=bf]{caption}
\usepackage[hidelinks]{hyperref}
\setlength{\parskip}{3pt}
\setlength{\textfloatsep}{8pt}
\setlength{\floatsep}{6pt}
\newcommand{\Qsq}{Q^2}

\title{\vspace{-1.2cm}ME52037 Task 2: Surrogate-Based Optimisation of a 50-Call Black Box}
\author{Oskar Chu}
\date{}

\begin{document}
\maketitle
\vspace{-1.2cm}

%==========================================================
\section{Problem and strategy}
%==========================================================
The black box returns a deterministic marketing score, MCSI ($0$--$10$), for three settings:
advert spend $x_1\in[0,200]$\,\pounds k, targeting spend $x_2\in[0,60]$\,\pounds k and audience skew $x_3\in[0,1]$, set to precisions of $0.1$, $0.1$ and $0.01$.
The aim was to maximise MCSI using at most 50 calls.
No gradient or shape information is available, so every call is expensive.
The strategy was therefore a \emph{surrogate-based} loop \cite{Forrester2008,Jones1998}: spend a fixed block of calls on a space-filling design, fit and validate a cheap model, then let the model choose each remaining call.
The budget was split as
\begin{equation*}
50 = \underbrace{25}_{\text{LHS training}} + \underbrace{5}_{\text{hold-out}} + \underbrace{17}_{\text{adaptive}} + \underbrace{3}_{\text{confirm}}.
\end{equation*}
All of this is implemented in a single MATLAB file, \texttt{bb\_session.m}, which logs every call, rebuilds the model after each reading, and can be resumed or undone.
Figure~\ref{fig:flow} shows the workflow.
The same code was rehearsed on synthetic ``fake boxes'' before the real session (Section~\ref{sec:verify}).

\begin{figure}[t]
\centering
\IfFileExists{figures/flowchart.png}{%
  \includegraphics[width=\textwidth]{figures/flowchart.png}}{%
  \fbox{\parbox[c][4.5cm][c]{0.95\textwidth}{\centering
  [Insert FigJam ``FINAL FLOW'' export here as \texttt{figures/flowchart.png}]}}}
\caption{Session workflow. Steps that use no box calls (planning, model building and the inner maximisation) are separated from the outer loop, which spends one call per iteration.}
\label{fig:flow}
\end{figure}

%==========================================================
\section{Initial sampling plan (calls 1--30)}
%==========================================================
\textbf{Size.}
The 10$d$ rule of thumb (Jones et al.~\cite{Jones1998}, examined by Loeppky et al.~\cite{Loeppky2009}) gives 30 points for $d=3$.
This was used as 25 training points plus 5 hold-out points, which leaves 20 calls (40\% of the budget) for adaptive search.
A sweep of the initial size on fake boxes (Figure~\ref{fig:ntrain}) supports this choice.
On a smooth single-peaked function, any $n_\text{train}$ between 10 and 30 ended within 0.1\% of the optimum.
At $n_\text{train}=35$ too few adaptive calls remained, and the search never switched to local refinement.

\textbf{Design.}
Inputs were scaled to the unit cube.
200 maximin Latin hypercubes \cite{McKay1979,Johnson1990} were generated with \texttt{lhsdesign}, each keeping the best of 50 random draws.
These were ranked jointly on minimum point-to-point distance (spread) and $L_2$-star discrepancy (uniformity) \cite{Warnock1972}.
The best plan was snapped to the setting precision.
LHS was preferred to a Halton sequence \cite{Kocis1997} because it guarantees one point per 1-D stratum in every input, so no setting range is left unsampled.
It also has no low-$n$ correlation artefacts, which unscrambled Halton points show in early dimensions.
The comparison over 100 draws per method (Table~\ref{tab:plan}) confirms this.
Scrambled Halton left a median of 8 empty 1-D slices, whereas every LHS variant left none.
The screened plan beats 97\% of single maximin LHS draws and all Halton draws on minimum distance.

\begin{table}[h]
\centering\small
\caption{Sampling plan quality, $n=25$, $d=3$ (medians of 100 draws; ``empty slices'' are unoccupied 1-D strata).}
\label{tab:plan}
\begin{tabular}{lcccc}
\toprule
Plan & Min.\ distance $\uparrow$ & $L_2$-star $\downarrow$ & $D^*$ $\downarrow$ & Empty slices \\
\midrule
Random & 0.084 & 0.0548 & 0.275 & 27 \\
Random LHS & 0.117 & 0.0354 & 0.176 & 0 \\
Maximin LHS (single) & 0.181 & 0.0346 & 0.173 & 0 \\
Scrambled Halton & 0.154 & 0.0307 & 0.158 & 8 \\
\textbf{Chosen (screened maximin LHS)} & \textbf{0.219} & \textbf{0.0263} & \textbf{0.138} & \textbf{0} \\
\bottomrule
\end{tabular}
\end{table}

\textbf{Hold-out points.}
Five extra points were placed greedily in the largest gaps of the design, which lie mostly near edges and corners.
They give an unbiased test of the model chosen on the 25 training points.
LOOCV alone is optimistically biased once it has been used to select the model \cite{Cawley2010}.
Because the gaps lie at the edges, the hold-out set acts as a worst-case extrapolation check.
After call 30 the hold-out points join the training data.

\begin{figure}[t]
\centering
\includegraphics[width=\textwidth]{figures/fig_ntrain_study.png}
\caption{Effect of initial sample size on an easy (top) and a deceptive two-peak (bottom) fake box, with 5 seeds per setting. (b) shows recommendation quality, (c) the call at which local refinement started and (d) the success rate. The chosen value $n_\text{train}=25$ is marked.}
\label{fig:ntrain}
\end{figure}

%==========================================================
\section{Surrogate modelling and validation (every call, 0 box calls)}
%==========================================================
Before every fit, points closer than the setting precision are merged.
Inputs are scaled to $[0,1]$ and MCSI is standardised.
Nine candidate models are then fitted:
\begin{itemize}
\setlength{\itemsep}{0pt}
\item \textbf{Polynomials} (linear, quadratic, cubic) via \texttt{fitlm}, with exact leave-one-out residuals from leverage (PRESS \cite{Allen1971}).
\item \textbf{RBFs} with a linear tail and five kernels: Gaussian, multiquadric, inverse multiquadric, cubic and thin-plate. Each width is tuned by Rippa's closed-form LOOCV \cite{Rippa1999}.
\item \textbf{Kriging} via \texttt{fitrgp}, with an ARD squared-exponential kernel, a constant mean and a fixed near-zero nugget, since the box is deterministic \cite{Sacks1989,Rasmussen2006}.
\end{itemize}
The model with the lowest LOOCV normalised RMSE is selected.
Its quality is reported as
\begin{equation}
\Qsq = 1-\frac{\sum_i \left(y_i-\hat y_{-i}\right)^2}{\sum_i \left(y_i-\bar y\right)^2}
     = 1-\text{NRMSE}^2\,\frac{n}{n-1},
\end{equation}
where $\hat y_{-i}$ is the prediction at point $i$ with that point removed.
$\Qsq$ is a predictive measure: unlike $R^2$, an interpolating model cannot score 1 on it trivially.
No single model family is best for every response \cite{Forrester2008}.
Selecting automatically at each call let the model change as data accumulated.
This happened in the session: inverse multiquadric RBF $\rightarrow$ Kriging $\rightarrow$ Gaussian and inverse multiquadric RBFs.

%==========================================================
\section{Inner level: maximising the surrogate (0 box calls)}
%==========================================================
The surrogate is cheap but can be multimodal, so it is maximised in three stages:
\begin{enumerate}
\setlength{\itemsep}{0pt}
\item A $21^3$ grid (9{,}261 points) is evaluated. The spacing of 0.05 in scaled units equals 10\,\pounds k, 3\,\pounds k and 0.05 skew, and the odd count includes the bounds.
\item Up to 6 of the best grid points are kept as seeds, each at least 0.25 apart. This places starts in different basins rather than six starts on the same peak.
\item \texttt{MultiStart} runs \texttt{fmincon} (SQP) from those seeds as a custom start-point set \cite{MathWorksMS}.
\end{enumerate}
SQP suits a smooth, bounded 3-variable problem: it converges superlinearly and respects bounds exactly \cite{Nocedal2006}.
Fixed seeds make the result deterministic.
Distinct optima more than 0.15 apart are kept: the best is the \emph{surrogate optimum}, and the others are \emph{rivals}, used by the outer level.
A sensitivity check (\texttt{maxcheck}) re-solved the final model with $11^3$--$41^3$ grids and 3--12 seeds.
Every combination returned the same optimum [130, 41.6, 0.30] to within $10^{-12}$ in predicted MCSI, in under 0.3\,s.
The default settings are therefore converged.

%==========================================================
\section{Outer level: choosing the next call}
%==========================================================
Each adaptive call is assigned one of four modes, checked in priority order:
\begin{enumerate}
\setlength{\itemsep}{0pt}
\item \textbf{Confirm} (3 or fewer calls left).
  The first two confirm calls step $+5\%$ of range in $x_1$ and then $x_2$ from the best measured point, to check that it is a true local maximum.
  The last confirm call measures the final surrogate optimum.
\item \textbf{Local} (once entered, the search stays local).
  The next point is the most isolated candidate inside the near-optimal box: the region within 2\% of range of the predicted maximum.
  Local mode is entered only when all of the checks in the box pass.

  \smallskip\noindent
  \fcolorbox{black}{gray!10}{\begin{minipage}{\dimexpr\linewidth-2\fboxsep-2\fboxrule\relax}
  \textbf{Decision checks for entering local mode (all must pass)}
  \begin{enumerate}
  \setlength{\itemsep}{0pt}
  \item \textbf{Stable:} adding the latest point moved the surrogate optimum by less than 0.02 in scaled units.
  \item \textbf{Accurate:} the last exploit or local call was predicted to within 1\% of the observed MCSI range.
  \item \textbf{No rival:} no second surrogate peak is predicted within 5\% of range of the best one.
  \item \textbf{Explored:} at least two explore calls have been made, and none beat the predicted optimum.
  \item \textbf{Budget:} at least 6 calls remain before the 3 confirm calls.
  \end{enumerate}
  \end{minipage}}
  \smallskip
\item \textbf{Exploit.} Query the surrogate optimum.
\item \textbf{Explore.} Query the most isolated point of 20{,}000 random candidates.
  The candidates are restricted to the top 25\% of predicted MCSI on alternate explore calls, and taken from anywhere in the domain on the others.
\end{enumerate}
Exploit and explore alternate in a cycle, as in Regis and Shoemaker's cycling of exploitation weights \cite{Regis2005}.
The cycle is exploit--exploit--explore when $\Qsq\geq0.8$ and no rival exists, otherwise exploit--explore.
A weak model therefore spends more of its calls exploring.
If the exploit point has already been sampled, the call becomes an explore.
The ``anywhere'' explore was added after rehearsals on a deceptive box, where a narrow global peak was hidden outside the high-prediction region.
After each call the actual settings, the prediction made before the call and the model code are logged.
The predict-then-query error is therefore a genuine out-of-sample test at every step.

\textbf{Recommendation.}
The final recommendation is the \emph{best measured} point, not the unmeasured surrogate optimum.
The box is deterministic, so a measured value carries no model error.
Measuring the surrogate optimum on the last call means the two coincide when the model is right.

%==========================================================
\section{Verification before the session}
\label{sec:verify}
%==========================================================
\begin{itemize}
\setlength{\itemsep}{0pt}
\item \textbf{Self-test.} The fast LOOCV formulas (PRESS for polynomials, Rippa for RBFs) were checked against brute-force refitting and matched to machine precision. Required toolboxes were confirmed present.
\item \textbf{Easy fake box} (smooth, single peak). Full rehearsals found the optimum within $<0.1\%$ for $n_\text{train}\le30$ (Figure~\ref{fig:ntrain}, top).
\item \textbf{Hard fake box.} This box has a broad rival peak of 8.58 and a narrow global peak of 9.00, both clipped to $[0,10]$ and rounded like the real box. Only 0--2 of 5 runs found the narrow peak (Figure~\ref{fig:ntrain}, bottom). This exposed a real limitation: a narrow peak lying between 25 samples is close to invisible. It motivated the explore-anywhere call and the measured recommendation, so the method at least never recommends a point worse than one already seen.
\end{itemize}

%==========================================================
\section{Results of the real session}
%==========================================================
\textbf{Fixed block.}
The best LHS reading was 2.46 (run 16).
At call 30, LOOCV on the 25 training points selected the multiquadric RBF, but with only $\Qsq=0.62$.
The hold-out test (Table~\ref{tab:holdout}) disagreed with that ranking.
The Gaussian RBF and Kriging generalised best to the edge points, while the multiquadric RBF did not.
This illustrates the selection bias the hold-out was designed to expose.
All five hold-out readings were $\approx0$, and every model under-predicted them, some below zero.
This is a sign of extrapolation at the domain edges rather than of an error in the interior.

\begin{table}[h]
\centering\small
\caption{Model check at call 30: trained on 25 LHS points, tested on the 5 hold-out points (NRMSE normalised by range).}
\label{tab:holdout}
\begin{tabular}{lccc}
\toprule
Model & LOO NRMSE & LOO $\Qsq$ & Hold-out NRMSE \\
\midrule
RBF multiquadric (selected) & 0.605 & 0.618 & 0.588 \\
RBF inverse multiquadric & 0.614 & 0.607 & 0.353 \\
RBF Gaussian & 0.649 & 0.562 & \textbf{0.283} \\
Kriging (ARD SE) & 0.737 & 0.434 & 0.294 \\
Polynomial quadratic & 0.770 & 0.382 & 0.854 \\
\bottomrule
\end{tabular}
\end{table}

\textbf{Adaptive calls.}
With $\Qsq<0.8$ the search started on the 1:1 cycle.
The first exploit (call 31) jumped to MCSI 3.98, against a prediction of 2.88, so the model under-predicted the peak by 27\% of range.
After refitting, Kriging was selected ($\Qsq=0.91$).
Its next exploit (call 32) read 4.03, within 0.5\% of its prediction.
Two explores followed, one in the top 25\% region (0.84) and one anywhere (0.31).
Neither found a better region.
With all decision checks passed, local refinement began at call 35.
During refinement, two manual explore overrides (calls 40 and 43) were spent re-testing distant regions for a missed peak; both returned below 0.9.
Local calls 35--47 all read 3.82--4.00.
Their prediction errors fell from about 2\% to below 0.2\% of range (Table~\ref{tab:calls}).

\begin{table}[h]
\centering\small
\caption{Adaptive and confirm calls. Error $=(y-\hat y)/\text{range}$, where $\hat y$ is predicted before the call. Models: IMQ/Gau = RBF inverse multiquadric/Gaussian, Krig = Kriging.}
\label{tab:calls}
\setlength{\tabcolsep}{4pt}
\begin{tabular}{rlccccrl@{\hspace{1.2em}}rlccccrl}
\toprule
Run & Mode & $x_1$ & $x_2$ & $x_3$ & $y$ & Err\% & Mdl &
Run & Mode & $x_1$ & $x_2$ & $x_3$ & $y$ & Err\% & Mdl\\
\midrule
31 & exploit & 134.3 & 43.1 & 0.26 & 3.981 & $+27.2$ & IMQ  & 41 & local   & 122.1 & 39.2 & 0.27 & 3.945 & $+1.3$ & Gau\\
32 & exploit & 130.3 & 42.5 & 0.28 & 4.030 & $+0.5$  & Krig & 42 & local   & 131.1 & 41.3 & 0.35 & 4.003 & $-0.1$ & IMQ\\
33 & explore & 149.4 & 58.5 & 0.00 & 0.840 & $-10.1$ & Krig & 43 & expl.*  & 198.9 & 22.2 & 0.87 & 0.060 & $-4.5$ & IMQ\\
34 & expl.\,any & 4.0 & 30.5 & 0.98 & 0.310 & $-6.0$ & Krig & 44 & local   & 120.4 & 44.6 & 0.25 & 3.867 & $+0.9$ & IMQ\\
35 & local   & 120.5 & 38.7 & 0.34 & 3.907 & $+1.8$  & Krig & 45 & local   & 132.8 & 38.3 & 0.30 & 3.974 & $-0.1$ & IMQ\\
36 & local   & 141.7 & 38.6 & 0.35 & 3.871 & $+0.5$  & Krig & 46 & local   & 140.8 & 41.7 & 0.31 & 3.976 & $-0.2$ & IMQ\\
37 & local   & 141.2 & 37.9 & 0.25 & 3.818 & $-0.3$  & Krig & 47 & local   & 119.5 & 42.1 & 0.30 & 3.981 & $+0.1$ & IMQ\\
38 & local   & 141.1 & 44.6 & 0.35 & 3.863 & $+0.7$  & Krig & 48 & conf.\,$+x_1$ & 140.3 & 42.5 & 0.28 & 3.964 & $0.0$ & IMQ\\
39 & local   & 122.1 & 44.6 & 0.35 & 3.907 & $+1.5$  & Krig & 49 & conf.\,$+x_2$ & 130.3 & 45.5 & 0.28 & 3.935 & $0.0$ & IMQ\\
40 & expl.*  & 196.2 & 57.3 & 0.33 & 0.841 & $-11.9$ & Gau  & \textbf{50} & \textbf{conf.\,opt} & \textbf{130.0} & \textbf{41.6} & \textbf{0.30} & \textbf{4.048} & $0.0$ & Gau\\
\bottomrule
\multicolumn{16}{l}{\footnotesize * manual explore override during local refinement.}
\end{tabular}
\end{table}

\textbf{Confirmation and recommendation.}
Stepping $+5\%$ of range from the best point at that time (call 32) lowered MCSI by 1.6\% in $x_1$ and 2.4\% in $x_2$.
This confirms that the best point lies on a local maximum, not a slope.
The final call measured the surrogate optimum at 4.0482, against a prediction of 4.0487, and became the best reading of the session.

\begin{center}
\fbox{\parbox{0.92\textwidth}{\centering
\textbf{Recommended setting:} $x_1=130.0$\,\pounds k, $x_2=41.6$\,\pounds k, $x_3=0.30$ \quad$\Rightarrow$\quad \textbf{measured MCSI = 4.0482} (run 50).\\
Worst case if each setting is off by one precision step: 4.0439 predicted ($-0.11\%$).}}
\end{center}

The final model is a Gaussian RBF with LOOCV $\Qsq=0.985$; the inverse multiquadric RBF (0.983) and Kriging (0.972) score almost the same, so the conclusion does not depend on the model choice.
The parity plot (Figure~\ref{fig:final}, left) is tight near the optimum, where 16 points read above 3.8 (within 6\% of the maximum).
The largest scatter is at low MCSI.
The slices (Figure~\ref{fig:final}, right) show a single interior peak.
Advert and targeting are optimal at about 65--70\% of range, and skew is optimal at about 30\%.
Skew is the most sensitive setting: MCSI falls to near zero at high skew.
Relative to the best LHS point, the adaptive phase improved MCSI by 64\%.

\begin{figure}[t]
\centering
\includegraphics[width=0.37\textwidth]{figures/fig_parity.png}\hfill
\includegraphics[width=0.61\textwidth]{figures/fig_final_model.png}
\caption{Final model (RBF Gaussian, $s=0.357$, $\Qsq=0.985$, 50 points). Left: LOOCV parity. Right: slices through the recommendation and 1-D sensitivities. White points are the fixed design, orange points are adaptive calls and the star is the recommendation. The black contour bounds the region within 2\% of range of the maximum.}
\label{fig:final}
\end{figure}

%==========================================================
\section{Discussion and limitations}
%==========================================================
\textbf{What worked.}
After the first exploit, predict-then-query errors at exploit and local calls stayed below 2\% of range.
The rigorous decision checks prevented premature local search: two explores were completed and failed to beat the optimum first.
In the end the measured and predicted optima agreed to $5\times10^{-4}$.
The model was robust to the choice of inner optimiser settings and nearly robust to the choice of model.

\textbf{Weak early model.}
At call 30, $\Qsq=0.62$ was low, and the hold-out check preferred different models.
The first exploit still found the right region, because even a poorly scaled model located the peak correctly.
Its magnitude was wrong by 27\% of range.
With hindsight, overriding the selected model at call 30 with the hold-out winner (Gaussian RBF or Kriging) would have been justified.

\textbf{Possible missed peak.}
The recommended MCSI of 4.05 is only 40\% of the 0--10 scale.
This may simply be the true maximum, but it cannot be proven from 50 samples.
The hard-box rehearsal showed that a narrow peak between samples can be missed.
Against that, four explores spread across the domain (calls 33, 34, 40 and 43) all read below 0.9, and every hold-out and LHS corner point was low.
Any undiscovered peak would therefore have to be narrower than the sample spacing of roughly 0.2 in scaled units.

\textbf{Modelling choices.}
The models allow negative predictions, as seen at the hold-out points, although MCSI is non-negative.
Modelling $\log(y+\epsilon)$ or clipping predictions would fix this, but it does not affect the optimum, which lies far from zero.
The local phase spent 13 calls refining a region already within 6\% of the optimum.
A trust-region step \cite{Forrester2008} or an expected-improvement criterion \cite{Jones1998} could have released some of those calls for further exploration.
Kriging was used even though the course notes cover RBFs.
It was included because it is the standard surrogate for deterministic computer experiments \cite{Sacks1989}, and LOOCV ranked it fairly against the others.

%==========================================================
\section{Conclusion}
%==========================================================
A screened maximin LHS, automatic LOOCV model selection with an independent hold-out check, a deterministic multistart inner optimiser and a rule-based exploit/explore/local/confirm outer loop together located a well-defined interior optimum within 50 calls.
The recommended setting $[130.0, 41.6, 0.30]$ gives a \emph{measured} MCSI of 4.0482.
It was confirmed as a local maximum by $\pm5\%$ steps and is insensitive to one-step setting errors.
The final surrogate is accurate ($\Qsq=0.985$).
The main residual risk, a narrow peak elsewhere, was tested by rehearsal and four domain-wide explores, and it is stated rather than assumed away.

\small
\begin{thebibliography}{99}
\setlength{\itemsep}{0pt}
\bibitem{Allen1971} D.~M.~Allen, ``The prediction sum of squares as a criterion for selecting predictor variables,'' Tech.\ Rep.\ 23, Dept.\ of Statistics, Univ.\ of Kentucky, 1971.
\bibitem{Cawley2010} G.~C.~Cawley and N.~L.~C.~Talbot, ``On over-fitting in model selection and subsequent selection bias in performance evaluation,'' \emph{J.\ Mach.\ Learn.\ Res.}, 11:2079--2107, 2010.
\bibitem{Forrester2008} A.~I.~J.~Forrester, A.~S\'obester and A.~J.~Keane, \emph{Engineering Design via Surrogate Modelling}, Wiley, 2008.
\bibitem{Johnson1990} M.~E.~Johnson, L.~M.~Moore and D.~Ylvisaker, ``Minimax and maximin distance designs,'' \emph{J.\ Stat.\ Plan.\ Inference}, 26:131--148, 1990.
\bibitem{Jones1998} D.~R.~Jones, M.~Schonlau and W.~J.~Welch, ``Efficient global optimization of expensive black-box functions,'' \emph{J.\ Global Optim.}, 13:455--492, 1998.
\bibitem{Kocis1997} L.~Kocis and W.~J.~Whiten, ``Computational investigations of low-discrepancy sequences,'' \emph{ACM Trans.\ Math.\ Softw.}, 23(2):266--294, 1997.
\bibitem{Loeppky2009} J.~L.~Loeppky, J.~Sacks and W.~J.~Welch, ``Choosing the sample size of a computer experiment: a practical guide,'' \emph{Technometrics}, 51(4):366--376, 2009.
\bibitem{MathWorksMS} The MathWorks, ``MultiStart'' and ``fmincon'', Global Optimization / Optimization Toolbox documentation, R2025b.
\bibitem{McKay1979} M.~D.~McKay, R.~J.~Beckman and W.~J.~Conover, ``A comparison of three methods for selecting values of input variables in the analysis of output from a computer code,'' \emph{Technometrics}, 21(2):239--245, 1979.
\bibitem{Nocedal2006} J.~Nocedal and S.~J.~Wright, \emph{Numerical Optimization}, 2nd ed., Springer, 2006.
\bibitem{Rasmussen2006} C.~E.~Rasmussen and C.~K.~I.~Williams, \emph{Gaussian Processes for Machine Learning}, MIT Press, 2006.
\bibitem{Regis2005} R.~G.~Regis and C.~A.~Shoemaker, ``Constrained global optimization of expensive black box functions using radial basis functions,'' \emph{J.\ Global Optim.}, 31:153--171, 2005.
\bibitem{Rippa1999} S.~Rippa, ``An algorithm for selecting a good value for the parameter $c$ in radial basis function interpolation,'' \emph{Adv.\ Comput.\ Math.}, 11:193--210, 1999.
\bibitem{Sacks1989} J.~Sacks, W.~J.~Welch, T.~J.~Mitchell and H.~P.~Wynn, ``Design and analysis of computer experiments,'' \emph{Statist.\ Sci.}, 4(4):409--435, 1989.
\bibitem{Warnock1972} T.~T.~Warnock, ``Computational investigations of low-discrepancy point sets,'' in \emph{Applications of Number Theory to Numerical Analysis}, Academic Press, 1972, pp.~319--343.
\end{thebibliography}

\end{document}
